% AI-produced by GPT-6 Astra (OpenAI), ChatGPT Pro mode, 18 November 2026.
% Working identifiers; editor assigns the continuing Part B numbers.
\subsection{Valuation inertia, centres, and boundary characters}

Here $K$ has characteristic zero, $L/K$ is finitely generated, and $v$
is a valuation of $L$ trivial on $K$.  Fix a prolongation to
$\overline L$; all decomposition and inertia groups are actual closed
subgroups of $G_L$.  Put $T=\widehat{\mathbb Z}(1)$ and
$\Gamma=\Gamma_v$.  We distinguish
\[
 \Gamma_{\rm div}=\bigcap_{n\geq1}n\Gamma
 \quad\hbox{from the largest divisible \emph{convex} subgroup}
 \quad\Delta\subset\Gamma.
\]
The first subgroup is divisible: if $a\in\Gamma_{\rm div}$, the
unique element $a/m$ in $\Gamma\otimes\mathbb Q$ belongs to
$n\Gamma$ for every $n$, since $a\in mn\Gamma$.  Conversely every
divisible subgroup is contained in it.  It need not be convex.

The page marks in the supplied second-pass transcription put scan~119
at manuscript page~59 and scan~121 at page~60; scan~120 continues the
former page and ends with the margin number~119.  Thus the question
in scan~120 belongs to manuscript page~59.  This is a check of the
supplied page-mark transcription, not a claim of a new inspection of an
unprovided image of that scan.

\begin{claim}{B.N3.1}{Proved}{The inertia formula in every rank}
There are canonical Kummer identifications
\[
 I_v=\operatorname{Hom}((\Gamma\otimes\mathbb Q)/\Gamma,\mu_\infty)
     =\operatorname{Hom}(\Gamma,T).
\]
The topology is the inverse-limit topology of finite Kummer quotients.
Furthermore
\[
 1\longrightarrow I_v\longrightarrow D_v\longrightarrow
 G_{\kappa(v)}\longrightarrow1
\]
splits, noncanonically.  Conjugation on inertia is the cyclotomic action.
If $r_\ell=\dim_{\mathbb F_\ell}\Gamma/\ell\Gamma$, then
\[
 I_v\simeq\prod_\ell\mathbb Z_\ell(1)^{r_\ell},\qquad
 r_\ell\leq\dim_{\mathbb Q}(\Gamma\otimes\mathbb Q)
 \leq\operatorname{trdeg}_K L-\operatorname{trdeg}_K\kappa(v).
\]
For $\gamma\in\Gamma$, evaluation
$\chi\mapsto\chi(\gamma)$ is identically zero exactly when
$\gamma\in\Gamma_{\rm div}$.  In particular $I_v=1$ exactly when
$\Gamma$ is divisible.  These statements re-prove FLAGGED\_149 D4
in the characteristic-zero setting at issue here.
\end{claim}
\begin{proof}
Let $H=L_v^h$ be the chosen henselization, and let $F=H^{\rm nr}$
be its maximal unramified extension.  Its residue field is
$\kappa(v)^{\rm sep}$ and its value group is $\Gamma$.  Units of
$F$ have $n$th roots in $F$: choose an $n$th root of the residue and
apply Hensel's lemma, whose derivative is invertible in residue
characteristic zero.  Tame extension theory, in the form of Temkin,
Theorem~3.4.4, and the absence of wild inertia in Corollary~3.5.3,
therefore identifies the absolute Galois group of $F$ with the stated
character group.  Explicitly its pairing is
\[
 (\sigma,v(a)/n\bmod\Gamma)\longmapsto
 \sigma(a^{1/n})/a^{1/n}.
\]
Independence of $a$ with a fixed value follows from the assertion about
units.  The finite Kummer pairings are perfect, since a value not in
$n\Gamma$ gives a nontrivial root extension.  Their union and dual
give the first formula.  Restriction to
$(n^{-1}\Gamma)/\Gamma\simeq\Gamma/n\Gamma$ gives the second.
Thus evaluation at $\gamma$ vanishes exactly when
$\gamma/n\bmod\Gamma=0$ for every $n$, which is the stated
divisibility criterion.  The exact decomposition sequence follows
from lifting finite separable residue extensions to unramified
extensions of the henselian valuation ring.

Here is a complement argument that does not assume a coefficient
field.  Among algebraic extensions $M/H$ in $\overline L$ with
unchanged residue field choose a maximal one.  Unions of chains
preserve this property.  If a value in $\Gamma_M$ is not divisible by
a prime $\ell$, adjoining an $\ell$th root of an element with that
value has ramification index $\ell$, hence unchanged residue field,
contrary to maximality.  Thus $\Gamma_M$ is divisible.  All these
extensions are henselian.  Residue-characteristic-zero henselian
fields are defectless (Temkin, Corollary~3.5.8), so $M^{\rm nr}$,
which has divisible value group and algebraically closed residue
field, is algebraically closed.  Hence restriction identifies
$G_M$ with $G_{\kappa(v)}$, and $G_M\hookrightarrow G_H=D_v$
is a complement.  Applying a decomposition element to the displayed
Kummer pairing fixes the value and changes only the root of unity;
this proves the cyclotomic assertion.

The torsion group $(\Gamma\otimes\mathbb Q)/\Gamma$ is divisible;
its $\ell$-torsion is $\Gamma/\ell\Gamma$.  The classification of
divisible torsion groups gives the displayed product after duality.
An $\mathbb F_\ell$-independent set of values is rationally
independent: in an integer relation divide by the greatest common
power of $\ell$ and reduce modulo $\ell$.  For the last bound,
choose elements with rationally independent values and units with
algebraically independent residues.  In a polynomial relation, terms
with distinct value exponent vectors have distinct values, and the
coefficient of the least-value vector is a nonzero polynomial in the
independent residues.  Cancellation is impossible.  These elements
are therefore algebraically independent, proving the bound.
\end{proof}

\begin{claim}{B.N3.2}{Proved}{The inertia normalizer and its necessary coarsening}
Suppose $I_v\ne1$.  Let $w=v/\Delta$ be the coarsening defined by
the largest divisible convex subgroup of $\Gamma$.  Then
\[
 I_w=I_v,\qquad N_{G_L}(I_v)=D_w.
\]
Writing $P_{L/K}=\operatorname{Gal}(\overline L/L\overline K)$ gives
\[
 N_{P_{L/K}}(I_v)=D_w\cap P_{L/K},\qquad
 1\to I_v\to D_w\cap P_{L/K}\to
 G_{\kappa(w)\overline K}\to1.
\]
Moreover $D_w=D_v$ exactly when the residual valuation $v/w$ on
$\kappa(w)$ is henselian.  In particular $N_{G_L}(I_v)=D_v$ for
every rank-one valuation with nontrivial inertia.  If $I_v=1$, its
normalizer is all of $G_L$.
\end{claim}
\begin{proof}
Convex subgroups form a chain, so their divisible members have a
largest union $\Delta$.  For any convex subgroup $A$, both $A$ and
$\Gamma/A$ are torsion free.  The exact sequence of divisible torsion
groups associated with $0\to A\to\Gamma\to\Gamma/A\to0$ and
Kummer duality gives
\[
 1\to I_{v/A}\to I_v\to\operatorname{Hom}(A,T)\to1.
\]
The first map is the actual inclusion for the induced prolongations:
its characters are pullbacks along $\Gamma\to\Gamma/A$.  It is an
equality exactly when $A$ is divisible.  Thus $I_w=I_v$.

On $F=\overline L^{I_v}$ the valuation $v$ is henselian, with
separably closed residue and value group $\Gamma$.  Its coarsening
by $A$ has separably closed residue exactly when $A$ is divisible.
Indeed its residual refinement is henselian, has value group $A$
and separably closed final residue; Claim~B.N3.1 computes the absolute
Galois group of the intermediate residue field as
$\operatorname{Hom}(A,T)$.  Consequently $w|_F$ is the coarsest
such coarsening.

It is also the coarsest among \emph{all} henselian valuations of $F$
with separably closed residue.  The precise comparison theorem used
here is Kuhlmann, Theorem~9.44 and Corollary~9.45: two independent
nontrivial henselian valuations force the field to be separably
closed, and two incomparable henselian valuations have a finest common
coarsening with separably closed residue.  If another such valuation
were incomparable with $w|_F$, their common coarsening would contradict
the preceding maximality of $\Delta$.  If it were a proper coarsening,
the same maximality would give a contradiction.  It must therefore
refine $w|_F$.  Since $I_v\ne1$, the field $F$ is not separably
closed, and this distinguished valuation is nontrivial.

An element normalizing $I_v$ acts on $F$ as an $L$-automorphism, so
preserves this intrinsically characterized valuation.  Its extension
to $\overline L$ is unique by henselianity, hence the element belongs
to $D_w$.  Conversely $D_w$ normalizes its inertia.  Intersection with
$P_{L/K}$ gives the geometric identity.  Constants have value zero;
passing from $L$ to $L\overline K$ preserves the value group and
adjoins $\overline K$ to the residue field.  The residue exact
sequence gives the last displayed sequence.  Finally $D_v$ is the
inverse image in $D_w$ of the decomposition group of $v/w$; that
group is the whole residue Galois group exactly for a henselian
valuation.  For rank one the only convex subgroups are $0,\Gamma$,
and $I_v\ne1$ excludes a divisible $\Gamma$.

This verifies FLAGGED\_149 D5, rather than replacing $\Delta$ by
the generally nonconvex subgroup $\Gamma_{\rm div}$.  It also checks
the two counterexamples in D6.  For the $x$-adic valuation of
$\mathbb Q(x,y)$, the geometric quotient of the normalizer is
$G_{\overline{\mathbb Q}(y)}\ne1$.  For an $x$-adic valuation of
$\mathbb Q(x,y,z)$ refined by the divisible-value valuation of
Claim~B.N3.3 on $\mathbb Q(y,z)$, the value group is
$\mathbb Z\oplus_{\rm lex}\mathbb Q$, so the two inertias agree.
The refinement is not henselian: the two prolongations to
$\mathbb Q(y,z)(\sqrt{1+y})$ distinguish residues $1,-1$ of the
root, while the quadratic polynomial is irreducible.  Thus its
decomposition subgroup, and hence $D_v$, is proper in $D_w$.
\end{proof}

\begin{claim}{B.N3.3}{Proved}{Divisible-value counterexamples, including nonzero invisible inertia}
FLAGGED\_149 D3 is valid: $K(x,y)$ has valuations with residue $K$
and value group $\mathbb Q$, or $\mathbb Z[1/2]$, centred at $(0,2)$
on $(\mathbb P^1)^2$.  Consequently the centre converse for arbitrary
valuations is false on $(\mathbb P^1-\{0,1,\infty\})^2$.

More strongly, on
$U=(\mathbb P^1-\{0,1,\infty\})^3$ there is a rank-one valuation
of $K(U)$ with $I_v\simeq T\ne1$, whose image in $\pi_1(U)$ is
trivial, but which has no centre on $U$.  For this example
$\Delta=0$; the inertia-preserving coarsening of Claim~B.N3.2 does
not repair the centre criterion.
\end{claim}
\begin{proof}
In the Hahn field $K((t^{\mathbb Q}))$ put
\[
 N_j=j!\quad\hbox{or}\quad N_j=2^j,\qquad
 a_j=j+1/N_j,\qquad x=t,\quad y=2+\sum_{j\ge1}t^{a_j}.
\]
The increasing supports are well ordered.  Their rational expressions
are series with support finite below any real bound, so sums,
products, and inverses obtained by the geometric-series expansion
are defined.  The value group is exactly
$\bigcup_m N_m^{-1}\mathbb Z$.  To prove the nontrivial inclusion,
let
\[
 P_m(X,Y)=\operatorname{Norm}_{K(X^{1/N_m})/K(X)}
 \left(Y-2-\sum_{j\le m}X^{a_j}\right).
\]
The $N_m$ conjugate partial sums are distinct: the numerator of the
last exponent is coprime to $N_m$.  At $(t,y)$ the identity factor
has value $a_{m+1}$; all other factors have values in
$N_m^{-1}\mathbb Z$.  Hence
$v(P_m(t,y))=a_{m+1}+c_m$ with $c_m$ in that group.  Starting from
$v(x)=1$ and $v(y-2)=a_1$, induction yields every $1/N_m$.
The reverse inclusion follows by examining supports.  Algebraicity
of $y$ over $K(t)$ would bound the value-group index over
$\mathbb Z$ by the finite field degree, contradicting these
unbounded denominators.  Thus $K(x,y)$ is a rational function field.
Leading coefficients lie in $K$, so the residue field is $K$.
The centre is $(0,2)$, on the single boundary branch $x=0$.
For the factorial choice $\Gamma=\mathbb Q$ and
Claim~B.N3.1 gives $I_v=1$.  The uniqueness of a centre on the
proper separated compactification rules out a centre on $U$.

For the stronger example choose a positive irrational real number
$\alpha$, and work in $K((t^{\mathbb Q\alpha+\mathbb Z}))$ with
the order inherited from $\mathbb R$.  Use the factorial sequence and
put
\[
 x=t^\alpha,\qquad
 y=2+\sum_{j\ge1}t^{\alpha a_j},\qquad z=2+t.
\]
The preceding norm argument gives value group $\mathbb Q\alpha$
on $K(x,y)$.  The element $t$ is transcendental over this field:
in a polynomial $\sum b_jt^j$, the nonzero terms have pairwise
distinct values $q_j\alpha+j$.  No cancellation is possible.
Therefore $x,y,z$ are algebraically independent, the total value
group is $\mathbb Q\alpha\oplus\mathbb Z$, and the residue is $K$.
Its largest divisible subgroup is $\mathbb Q\alpha$ (project to the
second summand), whereas its only convex subgroups are $0,\Gamma$
because it is archimedean.  Thus $\Delta=0$ and
$I_v=\operatorname{Hom}(\Gamma,T)=T$.
Its centre is $(0,2,2)$, with only $x=0$ on the boundary.  The local
inertia map is evaluation at $v(x)=\alpha$, which vanishes because
$\alpha\in\Gamma_{\rm div}$.  The other two curve coordinates have
interior centres.  The product formula for geometric fundamental
groups, or Claim~B.N3.7 below, now proves that its entire global
inertia image is trivial.  This is a nonzero inertia counterexample
and does not use a zero source group.
\end{proof}

\begin{claim}{B.N3.4}{Proved}{The local calculation and its exact image}
Let $X/K$ be smooth proper and $D=X\setminus U$ a simple normal
crossings divisor.  At a geometric centre $\bar x$ of $v$, write
$z_1,\ldots,z_r$ for the branch parameters and $\gamma_i=v(z_i)$.
The local geometric group is $T^r$ and the inertia homomorphism is
\[
 \rho_v:\operatorname{Hom}(\Gamma,T)\longrightarrow T^r,
 \qquad\chi\longmapsto(\chi(\gamma_1),\ldots,\chi(\gamma_r)).
\]
Its dual is
\[
 \psi_v:(\mathbb Q/\mathbb Z)^r\longrightarrow
 (\Gamma\otimes\mathbb Q)/\Gamma,
 \qquad(q_i)\longmapsto\sum q_i\gamma_i\bmod\Gamma.
\]
Its image is the annihilator of $\ker\psi_v$.  It is open exactly
when that kernel is finite, and surjective exactly when that kernel
vanishes.  Openness requires rational independence of the
$\gamma_i$; under independence its exact criterion is
\[
 [\Gamma\cap\operatorname{span}_{\mathbb Q}(\gamma_i):
       \textstyle\sum_i\mathbb Z\gamma_i]<\infty.
\]
Surjectivity means this index is one.  This re-proves
FLAGGED\_149 E3, including valuations of every rank.
\end{claim}
\begin{proof}
In the strict henselization of the local ring at $\bar x$, adjoining
all $n$th roots of the branch equations gives covers with group
$(\mu_n)^r$.  They are cofinal.  Indeed normalize in an arbitrary
finite cover of the punctured strictly local scheme and choose $n$
divisible by its finite tame inertia orders at the branches.
Abhyankar's lemma kills that codimension-one ramification after the
Kummer base change, whose target is regular.  Purity then makes the
normalization finite \emph{\'etale}, and strict henselianity makes it
trivial.  Normalization is finite here by excellence.  This proves
the local $T^r$ description without an assertion about its global
injectivity.

Pulling back $z_i^{1/n}$ to the valuation field gives the character
$\gamma_i/n\bmod\Gamma$ of Claim~B.N3.1.  Passing through every
$n$ proves both formulas.  Replacing $z_i$ by a unit multiple does
not change them, because units have all roots in the strictly
henselian valuation field.  Character duality is exact, since
$\mu_\infty$ is divisible and hence injective as an abelian group.
It identifies the quotient of $T^r$ by the compact image with the
dual of $\ker\psi_v$, giving the two image tests.

An integer relation among the values supplies an infinite copy of a
quotient of $\mathbb Q/\mathbb Z$ in the kernel.  Conversely, under
independence, the map $\mathbb Q^r\to\Gamma\otimes\mathbb Q$
is injective, so its preimage of $\Gamma$, modulo $\mathbb Z^r$,
is exactly the saturation quotient displayed above.  For a DVR
the formula reads $a\mapsto(m_1a,\ldots,m_ra)$ with $m_i>0$.
It is not open for $r>1$.  For the $\mathbb Z[1/2]$ example of
Claim~B.N3.3 on the one-branch stratum, the image is
$\prod_{\ell\ne2}\mathbb Z_\ell(1)$, also not open.  Thus
rational rank equal to the number of branches is not enough.
\end{proof}

\begin{claim}{B.N3.5}{False as stated}{Local boundary groups need not inject globally}
Even for a product of two hyperbolic curves, a smooth proper simple
normal crossings compactification can have a noninjective boundary
homomorphism $T^r\to\pi_1(U)$.  For
$U=(\mathbb P^1-\{0,1,\infty\})^2$, blow up $(0,2)$ in
$(\mathbb P^1)^2$.  At the intersection of the exceptional divisor
and the strict transform of $x=0$, the homomorphism is
\[
 T^2\longrightarrow T\hookrightarrow\pi_1(U),\qquad
 (a,b)\longmapsto a+b.
\]
Its kernel is the antidiagonal.  This counterexample concerns good
smooth SNC compactifications; it does not assert that the blow-up is
a compactification adapted as a smooth relative curve at every level
of the original tower.
\end{claim}
\begin{proof}
The chart $x=uv$, $y=2+v$ has boundary $uv=0$.  Each positive local
meridian goes once around $x=0$.  The $y$ coordinate lies in a
neighbourhood of the interior point $2$, so gives no meridian in the
second factor.  These identities hold in every finite Kummer
quotient, hence in $T$.  Cusp inertia embeds in the first punctured
line: its standard generators give a free profinite group of rank
two, and a meridian is a primitive generator.  The kernel is
therefore exactly as stated.  Equivalently the calculation can be
made with complex positive loops and Riemann existence.

The nonzero kernel does \emph{not} contradict the discrete-valuation
centre criterion.  A DVR centred at this crossing has positive
weights $m_1,m_2$, and its inertia maps by multiplication by
$m_1+m_2>0$, an injective endomorphism of $T$.  The supplied program
\texttt{n3\_boundary\_certificate.py} checks the coordinate matrix,
its integer kernel, and the stated finite Smith calculations; the
positive-weight assertion for all integers is the displayed formula.
\end{proof}

\begin{claim}{B.N3.6}{Proved}{A larger valid form of the centre margin}
Let $U=U_d\to\cdots\to U_0=\operatorname{Spec}K$ be a tower of
smooth geometrically connected curves, each with a smooth proper
relative compactification and finite \emph{\'etale} relative boundary.
Suppose no factor has type $(0,0)$ or $(0,1)$.  If
$\Gamma_{v,\rm div}=0$, then
\[
 I_v\longrightarrow\pi_1(U)\text{ is trivial}
 \quad\Longleftrightarrow\quad
 v\text{ has a centre on }U.
\]
This holds in every dimension and for every ordered rank.  In
particular it proves the DVR assertion of FLAGGED\_149 D7 and
extends it to all finitely generated value groups and to groups such
as $\mathbb Z[1/2]$.  The forward implication is false without
the value-group restriction, even when $I_v\ne1$, by
Claim~B.N3.3.  The reverse implication needs no type or value-group
restriction.
\end{claim}
\begin{proof}
We first supply the fibre-injectivity fact needed when elliptic and
$\mathbb G_m$ factors occur; using a hyperbolic-only theorem there
would omit these cases.  Over an algebraically closed field of
characteristic zero, every curve in the tower is an aspherical
complex surface after spreading out and embedding a field of
definition into $\mathbb C$.  Here ``surface'' means a real
two-manifold.  Compact curves of genus at least one and punctured
curves other than the projective line are aspherical.  A smooth proper
pair with disjoint boundary sections is a locally trivial topological
bundle, by local triviality of the proper smooth family together with
isotopy of its finite sections.  The homotopy sequence therefore gives
a discrete exact sequence of fundamental groups with injective fibre
and an aspherical total space, inductively along the tower.

We verify that the fibre remains injective after \emph{full}
profinite completion.  Recall that a group is good if its cohomology
with every finite module agrees with continuous cohomology of its
profinite completion.  Free groups of finite rank and oriented compact
surface groups of positive genus are good.  Here is the usual
effacement proof, included to specify this input.  After a finite
cover the module action is trivial.  A class in degree one is killed
by the kernel of its homomorphism to the finite module.  A class in
degree two for a compact surface is killed by a cover whose degree is
divisible by the exponent of the module; such covers exist in every
positive degree, and pullback multiplies the orientation class by the
degree.  Higher cohomology vanishes; for free groups degrees at least
two vanish.  The same argument holds in every finite cover.  This
effacement criterion implies goodness: the injections into finite
coinduced modules identify the two universal cohomological
$\delta$-functors, since continuous finite-module cocycles are
effaced by passage to a sufficiently small open subgroup.  Thus the
comparison is an isomorphism in every degree.

For clarity, the extension step is also proved.  Let
$1\to A\to B\to C\to1$ be one of the discrete sequences, with
$A$ a finitely generated surface or free group and with $C$ good,
as are its finite-index subgroups.  Choose any finite-index
characteristic subgroup $A_0\subset A$ and set $Q=A/A_0$.
The quotient $B/A_0$ is an extension of $C$ by finite $Q$.
On a finite-index subgroup $C_0$ the outer action on $Q$ is trivial.
The centralizer of $Q$ in its inverse image maps onto $C_0$, with
kernel $Z(Q)$, so is a central extension classified by
$H^2(C_0,Z(Q))$.  Goodness makes this class continuous.  Its cocycle
factors through a finite quotient, hence becomes zero on a further
finite-index subgroup $C_1$.  A splitting there gives a finite-index
subgroup of $B/A_0$ meeting $Q$ trivially.  Taking its normal core
produces a finite quotient of $B$ which is injective on $Q$.
Because $A$ is finitely generated, characteristic finite-index
subgroups are cofinal.  Hence the induced topology on $A$ is its
full profinite topology, and
\[
 1\to\widehat A\to\widehat B\to\widehat C\to1
\]
is exact.  The Lyndon--Hochschild--Serre spectral sequences for the
discrete and profinite extensions now prove goodness of $B$:
$H^q(A,M)$ is finite for finite $M$, and goodness of $A$ and $C$
identifies their entire $E_2$ pages.  Finite-index subgroups have
the same extension structure and satisfy the same argument.  Induction
therefore proves the required injectivity at every stage.  Riemann
existence and invariance under algebraically closed characteristic-zero
extension transfer it to the geometric \emph{\'etale} groups.
For hyperbolic-only towers the same result is Hoshi (2014),
Proposition~2.4(i)--(ii); its published hypotheses are connected
noetherian separated normal bases over a characteristic-zero field.

If $v$ has a centre, pullback of every finite \emph{\'etale} cover
to the strict henselization of its valuation ring is constant.
Consequently inertia acts trivially.  Conversely suppose the image
of $I_v$ is trivial and extend the generic map successively along
the tower.  It extends at level zero.  Suppose it has extended to
$\operatorname{Spec}V\to U_{i-1}$.  The proper relative
compactification of $U_i$ extends the generic lift uniquely to
$\operatorname{Spec}V$.  If its closed point meets the relative
boundary, then after strict henselization a single boundary section
has local equation $z$ and $\gamma=v(z)>0$.  The induced inertia
map to its cusp group is evaluation at $\gamma$, by the Kummer
calculation in Claim~B.N3.4.  It is nonzero because
$\Gamma_{\rm div}=0$.  This cusp group embeds in the geometric
fibre group, which embeds in $\pi_1(U_i)$ by the preceding argument,
contradicting the assumed trivial image under
$\pi_1(U)\to\pi_1(U_i)$.

Here is a local-ring verification of that comparison which avoids
any noetherian hypothesis on $V$, and does not apply a DVR theorem
to it.  Let $b$ be the centre in $U_{i-1}$ and $x$ its lift on the
relative compactification.  The total space is smooth over $K$ and
therefore noetherian and regular.  Its strictly henselian local ring
$R=\mathcal O^{\rm sh}_{\overline U_i,\bar x}$ maps to
$V^{\rm sh}$ by the universal property of strict henselization,
after fixing the compatible separably closed residue fields.  Near
$x$ the relative boundary has one branch, so
$\pi_1(\operatorname{Spec}R[1/z])=T$ by the noetherian local
Kummer--purity argument of Claim~B.N3.4.  The punctured strictly
local ring in the geometric fibre at $\bar b$ also has group $T$;
its map to the total punctured local scheme induces the identity on
these $T$'s, since in both cases their finite characters are the
roots of the same $z$.  Thus the total local $T\to\pi_1(U_i)$
is exactly the cusp map in that geometric fibre followed by its
already proved injection.  Pullback along $R\to V^{\rm sh}$
computes it by evaluation at $v(z)$ using Claim~B.N3.1.
Strict henselization preserves the value group.  This proves the
comparison for every valuation ring directly through noetherian
local rings on the given tower.

For a proper elliptic factor there is no boundary, so properness
already extends the map.  For a $\mathbb G_m$ factor its cusp group
is $T$ and the same Kummer calculation applies.  For every hyperbolic
affine factor cusp inertia is injective.  One direct verification of
the latter is to pass to a finite cover with at least two punctures;
then the cusp has a nonzero primitive class in the free abelianization.
Such a cover exists for a once-punctured positive-genus curve by any
nontrivial unramified cover of its proper completion; genus-zero
hyperbolic curves already have at least three punctures.  A
nontrivial multiple of a cusp generator is still detected.  This
checks every allowed type.  Avoiding the boundary at the closed
point means avoiding it on all of the local scheme $\operatorname{Spec}V$.
The induction gives the required centre without a descent gap.
\end{proof}

\begin{claim}{B.N3.7}{Proved}{An exact answer for products, with every factor type}
Let $U=\prod_{j=1}^d C_j$ be a product of smooth geometrically
connected curves in characteristic zero, with proper smooth
completions.  Let $x$ be the centre of $v$ on their product, and for
each coordinate that is a missing point choose its local parameter
$z_j$ and write $\gamma_j=v(z_j)>0$.  Then
\[
 I_v\to\pi_1(U)\text{ is trivial}
 \quad\Longleftrightarrow\quad
 \gamma_j\in\Gamma_{\rm div}\quad(j\in B_{\ne\mathbb A^1}).
\]
Here $B_{\ne\mathbb A^1}$ is the set of boundary coordinates whose
curve factor is not geometrically isomorphic to $\mathbb A^1$.
Proper factors contribute no condition; affine-line factors
contribute no condition.  In particular projective-line and elliptic
proper factors, and punctured-line factors of type $(0,2)$, are all
included.
\end{claim}
\begin{proof}
Inertia is geometric because the valuation is trivial on $K$.
The geometric fundamental group of the product is the product of
the curve groups; this follows over $\mathbb C$ from the product
of topological groups and finite generation, and then by Riemann
existence and base-extension invariance.  At an interior coordinate
the inertia image is zero by the centre implication.  At a boundary
coordinate the image is evaluation at $\gamma_j$ followed by that
curve's cusp homomorphism.  This homomorphism is zero for
$\mathbb A^1$, whose geometric group is trivial, and is injective
for every other nonproper curve by the cusp proof in
Claim~B.N3.6 (including $\mathbb G_m$).  Independent factors cannot
cancel in a direct product.  Claim~B.N3.1 gives precisely the test
displayed above.  The argument is unchanged after the finite residue
extension needed to split a boundary point, since that extension is
unramified for the valuation.
\end{proof}

\begin{claim}{B.N3.8}{Proved}{The excluded factor types can kill a later meridian}
For $\mathbb A^1$ a DVR at infinity has no centre and trivial
geometric inertia image.  The type $(0,0)$ can also cause failure in
a nontrivial tower: the $\mathbb G_m$-bundle
\[
 U=\mathcal O_{\mathbb P^1}(-1)^\times\longrightarrow\mathbb P^1
\]
is an elementary tower of types $(0,0),(0,2)$, and its geometric
fundamental group is trivial.  A divisorial valuation along the zero
section of $\mathcal O(-1)$ has nonzero inertia and no centre on
$U$, yet has trivial global image.
\end{claim}
\begin{proof}
The first assertion follows from
$\pi_1(\mathbb A^1_{\overline K})=1$ in characteristic zero.
The map sending a nonzero vector in a line in $K^2$ to that vector
identifies the second total space with $\mathbb A^2\setminus\{0\}$.
Over $\mathbb C$ it retracts onto $S^3$, so is simply connected;
Riemann existence gives the geometric assertion over every
characteristic-zero field.  The zero section is a smooth divisor in
the line bundle and supplies a DVR with inertia $T$.  It has no
centre on the complement.  A good relative compactification of the
bundle is $\mathbb P(\mathcal O\oplus\mathcal O(-1))$, with its
zero and infinity sections as boundary.  In topological terms the
boundary map from $\pi_2(\mathbb P^1)$ to the fibre $\pi_1$ is
multiplication by $-1$, the degree of the line bundle.  This also
explains why a fibre-injection argument cannot include projective-line
bases indiscriminately.
\end{proof}

\begin{claim}{B.N3.9}{Proved}{A global criterion using finite-cover residues}
Keep the SNC compactification and centre of Claim~B.N3.4.  Let
$P_x$ be its punctured strictly local neighbourhood over
$\overline K$, with group $T^r$.  Define a subgroup
$\mathcal R_x\subset(\mathbb Q/\mathbb Z)^r$ as follows.
For every connected finite \emph{\'etale} cover
$W\to U_{\overline K}$ admitting a pointed lift
$P_x\to W$, for every $n>0$, and every class
$\alpha\in H^1(W,\mu_n)$, include its restriction to $P_x$.
The latter is the residue vector
$(a_1/n,\ldots,a_r/n)$ in the indicated group.  Then
\[
 I_v\to\pi_1(U)\text{ is trivial}
 \quad\Longleftrightarrow\quad
 \sum_i a_i\gamma_i\in n\Gamma
 \quad\text{for every such finite-cover residue vector}.
\]
This is valid for every smooth $U$, every elementary tower and every
factor type, and every valuation.  It uses cyclic cohomology on
finite covers, rather than an unspecified kernel of
$T^r\to\pi_1(U)$.
\end{claim}
\begin{proof}
Write $G=\pi_1(U_{\overline K})$ and $f:T^r\to G$ for the local
homomorphism.  A pointed lift to the cover defined by an open subgroup
$H\subset G$ exists exactly when $f(T^r)\subset H$; it is the
chosen component on which the local cover is trivial.  The residue
classes above are therefore precisely restrictions of finite-order
characters of these open subgroups.  They separate every nonidentity
element in the compact abelian subgroup $f(T^r)$.  To verify this,
let $a\ne1$ belong to that subgroup, choose a finite quotient
$q:G\to Q$ with $q(a)\ne1$, and put
$A=q(f(T^r))$ and $H=q^{-1}(A)$.  The finite group $A$ is abelian.
A cyclic character of $A$ detects $q(a)$; pulled back to $H$ it is
an element of $H^1(W,\mu_n)$ for the corresponding cover $W$.
It therefore belongs to the defining family and detects $a$.

It follows that the global inertia image is zero exactly when all
these characters vanish on the local image of $I_v$.  The dual
formula in Claim~B.N3.4 makes this condition
$\sum(a_i/n)\gamma_i=0\bmod\Gamma$, which is the displayed
divisibility test.  All parts of the construction are geometric
finite-cover and Kummer-cohomology data, and the valuation enters
only through finitely many values in each test.  The residue vector
can equivalently be computed by extending a cyclic torsor to the
normalization and reading its Kummer orders on the branches.
The proof uses arbitrary finite quotients, so it is a profinite
statement, not a bounded-quotient experiment.

Here is a description of the same module in terms of finite
\'etale algebras and divisors, specifying how to calculate its
characters.  For each such pointed cover $W$, put
$F_W=\overline K(W)$ and
\[
 \mathcal K_n(W)=
 \{h\in F_W^\times:\operatorname{div}_W(h)
       \in n\operatorname{Div}(W)\}/F_W^{\times n}.
\]
The Kummer identification $H^1(W,\mu_n)=\mathcal K_n(W)$
holds because $W$ is regular and $n$ is invertible.  Indeed, a
torsor restricts to $Y^n=h$ at the generic point;
unramifiedness at codimension one makes each divisor order a
multiple of $n$.  Conversely these divisibilities make the
normalization in $Y^n=h$ unramified at codimension one: locally
divide $Y$ by the appropriate uniformizer power, leaving the
\'etale equation $Y^n=u$ for a unit.  Purity on the regular $W$
makes the normalization \'etale everywhere.  Generic restriction
is injective because a generic section of a finite \'etale torsor
over a normal scheme extends uniquely.  This proves the
identification, including torsors from nontrivial $n$-torsion
line bundles; it is not restricted to global units.

The pointed lift embeds $F_W$ in the fraction field of the chosen
strictly local branch.  For $h\in\mathcal K_n(W)$, factor its
pullback in that regular strictly local ring, allowing poles.
All irreducible factors except the boundary parameters $z_i$
have orders divisible by $n$.  Units have $n$th roots by
Hensel's lemma.  Its character vector is therefore
\[
 \left(\frac{\operatorname{ord}_{z_1}(h)}n,\ldots,
       \frac{\operatorname{ord}_{z_r}(h)}n\right)
       \pmod{\mathbb Z^r}.
\]
Thus $\mathcal R_x$ may be defined using these divisor orders
alone, ranging over finite \'etale algebras with the stated local
section.  This definition does not require a previously unknown
kernel of $T^r\to G$.  For any specified cover and $n$, the input
is the divisor group, the rational functions whose divisors are
divisible by $n$, and the chosen local branch.

For the original product compactification of Claim~B.N3.7 the module
$\mathcal R_x$ is the full coordinate module on the non-affine-line
boundary factors.  For the blow-up in Claim~B.N3.5 it is the
diagonal subgroup $\{(q,q):q\in\mathbb Q/\mathbb Z\}$.
Thus its exact test there is
$\gamma_1+\gamma_2\in\Gamma_{\rm div}$, not separate
divisibility of the two values.  These computations exhibit how a
noninjective boundary group is handled without mistaking its local
image for its global one.
\end{proof}

\paragraph{Scope of the determination.}
Claim~B.N3.9 is an exact finite-cover cohomological criterion, with a
proof that the indicated abelian tests suffice even when the ambient
fundamental group is nonabelian.  It is not a finite algorithm from
the list of factor types, nor a computed classification of
$\mathcal R_x$ for every nonproduct tower.  This remains an
uncompleted part of N3, which is kept open in the task register.
The proofs above settle
the universal local-injectivity assertion negatively, verify D3--D5
and E3, verify D7 for DVRs with every allowed factor type, give the
larger value-group range of the centre margin, and give nonzero
invisible inertia outside that range.  They do not infer a global
centre theorem from local injectivity at a crossing.

\paragraph{Checks and provenance.}
The displayed arguments were written and reread by GPT-6 Astra
(OpenAI), ChatGPT Pro mode, on 18 November 2026.  There is no claim
of human or different-model verification.  The finite code checks
only the stated integral matrices and finite quotients.  Temkin's
valuation notes, Kuhlmann's Chapter~9, Hoshi's published paper,
the supplied SGA~1 source, and Stacks Tag~0BMB were actually opened.
No priority claim is made for these formulations.
